\section{System Integration and Verification}
\label{section:integrate-and-verify-system}

The following procedures specify future work. No empirical framework V\&V has run. OOSEM specifies verification subjects, stimuli, conditions, expected responses, verdicts and the enabling verification environment \citep[Secs.~17.3.8--17.3.9]{friedenthal2014practical}. NASA distinguishes requirement conformity from suitability for stakeholder expectations and ConOps. Both can be assessed through inspection, analysis, demonstration or test, with methods suited to product maturity \citep[Secs.~2.4, 5.3--5.4 and Appendix~I]{nasa2016systems}.

The assessment subjects remain distinct. Framework verification checks the implemented framework's FRs. Framework validation checks engineering suitability in representative use. EPS phase-product V\&V checks the engineered concept, requirements, models or host software. Model credibility assesses supported predictions, with calibration kept separate from independent assessment. Research evaluation measures the enterprise MOEs and bounded transfer. Shared observations can support several assessments only when each subject and comparison basis remain explicit.

Qualification assesses an identified future framework configuration or assistance profile. It uses approved criteria within the model and software scope. Hardware and flight qualification remain outside scope.

\subsection{Framework requirement verification}

The subject is an identified implemented framework candidate and core revision. Record component, authoring, rule, corpus, scoring-model, solver, compiler and library versions. Declare the host, workload, access rights, selected analysis tools, F Prime integration profile, reference data, backup and reconstruction arrangements. Verification personnel need the relevant systems, EPS, analysis and software skills. Independent experts establish reference results and adjudicate disagreements separately from development and blind assessment. Ordinary computing, controlled models/data and reviewers form the enabling environment. No physical EPS equipment or testing facility is required.

The Systems Engineer owns this plan and its criteria. This engineer agrees reference answers, intended-use limits, workloads, thresholds and comparison rules with engineers, independent experts and stakeholders at the framework requirements baseline, before qualification. Missing criteria block performance qualification. The designated authority controls baseline and procedure changes. Numerical targets remain to be resolved (TBR). Their risk basis includes missed syntax defects, unsuitable recommendations, incorrect numerical results and incomplete analysis reconstruction. Resolution covers recommendation applicability, change impacts, numerical errors, reconstruction and view interpretation. Fix workloads, thresholds and comparison rules before execution. Recording a TBR supplies no passing verdict \citep[Secs.~6.1.1.2.4 and~6.7]{nasa2016systems}.

\paragraph{Common procedure for every V case.}
(1) Identify the subject, applicable FR, configuration and controlled references. (2) Confirm prerequisites, profile coverage and approved criteria. Record missing inputs or unresolved criteria. (3) Apply the specified stimuli and conditions, including clean and failed controls. (4) Preserve raw responses and compare them with independent expected results. (5) Record each criterion, discrepancy and applicability limit against that subject. (6) Authorize corrections prospectively, retain predecessors and repeat affected checks. A changed implementation, profile or corpus requires impact assessment and repeated applicable qualification. An anomaly may originate in the product, procedure or environment. Resolve that distinction before assigning conformity.

Inspection (I) examines records and configurations. Analysis (A) compares independent calculations or labeled results. Demonstration (D) shows a capability. Test (T) obtains detailed implementation responses under controlled conditions. The 23 active cases, V01--V05, V07, V09--V11, V13--V17, V19--V20, V22, V25--V27, V34, V37 and V39, target the same-numbered FR. Table~\ref{tab:framework-verification-increments} and the common procedure specify the planned methods, prerequisites, stimuli, expected responses, comparisons and retained evidence.

{\small
\setlength{\tabcolsep}{4pt}
\setlength{\parskip}{0pt}
\renewcommand{\arraystretch}{1.06}
\begin{longtable}{@{}>{\raggedright\arraybackslash}p{0.24\linewidth}>{\raggedright\arraybackslash}p{\dimexpr0.76\linewidth-2\tabcolsep\relax}@{}}
\caption{Planned implementation verification groups and methods. Apply the common procedure to each case.}\label{tab:framework-verification-increments}\\
\toprule
\textbf{Cases / methods / prerequisites} & \textbf{Stimulus, expected response, comparison and retained evidence} \\
\midrule
\endfirsthead
\toprule \textbf{Cases / methods / prerequisites} & \textbf{Stimulus, expected response, comparison and retained evidence} \\
\midrule
\endhead
\bottomrule
\endfoot
V01, V25 (I,T)\par V04 (T), V17 (A,T)\par Approved requirement, mission and derivation references & Round-trip quantitative, qualitative, unresolved and inapplicable records. Query both trace directions, revise degraded modes and derive component allocations from sourced quantities. Compare exact fields, revisions and endpoint sets, plus independent derivation calculations under fixed numerical rules. Missing or ambiguous links and criteria remain explicit. Retain inputs, calculations, records, graph differences and engineer dispositions. \\
V02, V03 (A,T)\par V39 (A,D)\par Independent syntax, method and activity references & Edit clean drafts and drafts with seeded syntax errors. Vary requirements, available inputs and the current activity. Compare affected-text findings, recommended methods/models, proposed activities and unmet prerequisites with independent references. Retain inputs, responses, missed defects, unsuitable recommendations and comparison results. Qualification below supplies the statistical comparison. \\
V13, V14, V20 (A,T)\par V15 (I,D)\par Sourced technical properties, common mission and fixed trade criteria & Substitute components and configure EPS representations. Recalculate mass, volume, mode power, sunlight/eclipse energy and beginning/end-of-life array/battery sizing. Exercise discrete cell configurations, lifetime margins and electrical architecture choices. Include unknown properties and an infeasible eclipse case. Compare independent arithmetic, allocation margins and feasibility classifications. Preserve unknown values in the comparisons. Retain model/extension revisions, sources, equations, raw measures and differences. Numerical rules require prior agreement. \\
V16, V27 (A,T)\par V26 (I,T)\par Independent EPS and RAM references & Exercise each FR-16 analysis family with independent energy, photovoltaic, converter, battery, thermal and protection-stress references. Apply time-dependent loads, mode transitions and supplied environmental conditions. Compare electrical responses, dissipation, energy shortages and operating-limit margins. Identify applicable FR-26 analyses and execute the FR-27 representations on controlled component failure dependencies. Compare failure effects, criticality, minimal failure combinations, conditional probabilities and mission reliability with independent references. Include transient shortages, failed analyses and incomplete failure data. Retain inputs, model choices, raw outputs, comparisons and execution status. \\
V22 (A,T)\par Known model parameters and quality-checked supplied data & Estimate photovoltaic, battery electrical/electrothermal and dynamic EPS/load parameters from controlled datasets. Compare estimated parameters and model residuals with independent calculations under the same fitting method and inputs. Retain supplied data, model definitions, reference values, estimates and differences. Parameter fit supplies no independent evidence of physical prediction validity. \\
V19, V34 (D,T)\par Declared F Prime profile and independent MIL sequences & Check EPS component and port mappings into the declared F Prime profile, including version compatibility. Execute MIL scenarios for nominal operation, direct energy transfer, MPPT/end-of-charge transitions, command/mode changes, overcurrent/undervoltage/overvoltage faults, load shedding and recovery. Compare mappings and simulated state responses with independent references. Retain specifications, profile versions, mapped definitions, scenario inputs, observed responses and discrepancies. \\
V05 (A,T)\par V11 (T)\par V07, V09 (I,T)\par V10 (D,T)\par Identified baselines, reference impacts and analysis reconstruction inputs & Propose mission, requirement, component and model changes. Compare identified requirement and EPS design impacts with independent references. Retrieve exact baseline revisions and authorized decisions. Check result links and retained failed runs. Seed missing, rejected, stale, failed and unresolved evidence and compare the affected requirement reports. Independently reconstruct analyses from the supplied inputs, dependencies and instructions. Retain source records, retrieved baselines, reconstructed outputs and discrepancies. \\
V37 (I,D)\par Controlled requirement and analysis records and independent view references & Generate filtered mission, requirement-allocation, design, interface and analysis views. Compare selected records and displayed values with the source records. Ask fixed engineering interpretation questions and compare answers with independent references. Retain filters, generated views, answers and discrepancies. \\
\end{longtable}
}

Each future verification report retains the exact requirement/reference/subject/environment revisions, as-run procedure, raw data, comparison calculations, criterion outcomes, limitations, deviations and discrepancy dispositions. Pass, fail or undetermined describes only the named subject under the fixed rule. Planned cases have no observed verdict. Reduced mandatory profile scope requires a reviewed requirement/scope change before evaluation. A documented development gap cannot replace a required successful exchange.

\subsection{Active assistance qualification}

Use an independently labeled assessment set separate from rule and model development. Include clean requirements and each declared syntax defect family. Compare identified affected text with the reference labels. Report TP/FP/FN, precision and recall, missed defects and false findings.

Assess analysis, verification and validation recommendations against independent expert judgments for the supplied requirements and available inputs. Assess proposed engineering activities and identified prerequisites against references for the current activity. Include cases with incomplete inputs and no applicable method. Record recommendation errors and expert disagreement separately from syntax findings.

Qualification applies to the identified rule and recommendation configuration. Changes to that configuration require a justified comparison with the qualification basis.

\subsection{Framework validation in engineering use}

Anticipated Systems, EPS and Non-EPS Engineers perform representative lifecycle tasks, with mission stakeholders and independent reviewers. Include novice and experienced contributors. Declare training, staffing, available time, computing resources, data completeness, tool rights, collaboration and service availability representative of the team. Observe complete concept, requirements, preliminary design, detailed model/software, change and review tasks. Passing FR checks alone does not establish suitability \citep[Sec.~5.4 and Appendix~E]{nasa2016systems}.

Before execution, agree stakeholder expectations and engineering ConOps, expected phase products, task limits and success criteria. Independent judges prepare reference outcomes and fixed review questions. Engineers establish feasible concepts and refine and allocate requirements. They justify sizing and trades and assess declared analyses and MIL scenarios. They handle a coupled change and prepare and review a coherent successor. Include missing data, misleading assistance, failed analysis, stale evidence and clean controls. Judges assess outcome correctness, critical omissions, interpretation, reconstruction and whether the workflow meets actual needs under those conditions. Preserve products, task observations, decisions, support used, failures and unmet needs. Stakeholders determine acceptance against their expectations. That decision remains separate from expert classifier labels and requirement verdicts.

\subsection{EPS phase products and model credibility}

EPS V\&V can begin in Pre-Phase A with concept descriptions, budgets and models. Verification checks applicable EPS requirements. Validation assesses agreed operational energy/service and mode/recovery expectations from ConOps. Representative orbit, attitude, illumination, eclipse, load transitions, temperature and nominal/degraded conditions apply within the declared fidelity. Early analysis can validate mission suitability. MIL exercises modeled control with a simulated plant. SIL can assess released host code against simulated EPS/subsystem products as part of EPS product assessment. It is separate from the framework's retained MIL execution requirement. Each plan identifies the exact subject, requirement or expectation revision, inputs, conditions, expected response, criteria, retained observations and limitations \citep[Secs.~2.4, 5.3--5.4 and~6.1.1.2.4]{nasa2016systems}.

Independent equations, interface inspections, state/event references and supported histories provide comparison bases. Zhu motivates separate physics outputs streamed to scene/data entities, without demonstrating CUBICS, F Prime or bidirectional solver compatibility \citep[pp.~4--5, 7--8]{handpicked-0001}. Its SIL definition is separate from the dataset calibration/prediction assessment, which documents no executed SIL procedure \citep[p.~3 and pp.~7--9]{handpicked-0001}. CDR (MR-05) assesses detailed design during adapted C before deliverable coding. Exact current design/interface/software-scope authorization then controls coding. SIR (MR-06) assesses released products and scoped integration readiness at the end of C. These decisions remain distinct. Target qualification, physical prototypes/testing and PIL/HIL are scope exclusions.

Model credibility identifies intended prediction, fidelity, omitted mechanisms and supported domain. Check equations, units, implementation, coupling, convergence and temporal resolution against independent references. Calibrate on quality-checked supplied data, retaining prior/revised parameters and residual units. Assess prediction separately on withheld independent data across relevant load, temperature and health ranges where available. Missing independent data or error limits remain gaps. A good fit or repeatable run establishes neither predictive validity nor EPS suitability.

\subsection{Research evaluation and bounded reuse}

Fix two contrasting EPS cases before execution. One can use sparse data and a simple concept. The other can use variable loads, richer supplied data and detailed host control. Freeze the common reusable core, allowed mission parameters, variants/extensions and model/software mappings. Fix tasks, independent reference outcomes, training/participant assignment, baseline workflow, workload/time resources, comparison rules and declared technical resource limits. This protocol creates no operational requirement to complete two cases.

Both cases cover concept, requirement, design, software, change and review work with clean and seeded failed controls. Compare MOE-01--04 only for correct outcomes under comparable conditions. Report active person-hours separately from elapsed, waiting and unattended hours. Report early expert-confirmed issues, critical misses, unnecessary investigations, coverage, interpretation and reconstruction. Report learning/adaptation and declared computing-resource use separately. Preserve per-case results, uncertainty and unsuccessful outcomes.

Core changes require a prospective reviewed successor revision. Preserve original results, failures and adaptations. Rerun both cases on the changed core and disclose both reruns before making a bounded transfer claim. Do not hide mission-specific core edits as permitted parameters. The present chapter specifies design and planned assessment. Framework qualification requires an implemented configuration, agreed numerical criteria and empirical V\&V/evaluation.
